\section{Results --- Layer 2: Live Case Study (P6-B)}\label{results-layer-2-live-case-study-p6-b}

Layer 2 supplies what L1 cannot: an end-to-end, deterministic, \emph{ecologically-real} trace of the harness in action on real GitHub issues. The case study is consciously \emph{not} an LLM-generated transcript: it is a typed, reproducible flow through \passthrough{\lstinline!context\_insert → sublate\_with\_evidence → compact → final answer!} whose output is exactly the qualifier-set of the surviving live items. This isolates the \emph{context discipline} from generation quality.

\subsection{Cases}\label{cases}

We hand-picked three real issues drawn from popular Python projects, each with a documented mix of stale-vs-fresh evidence in the actual user-facing search results:

\begin{enumerate}
\def\labelenumi{\arabic{enumi}.}
\item
  \textbf{\passthrough{\lstinline!django\_request\_body!}} --- ``\emph{When does Django raise \passthrough{\lstinline!RawPostDataException!} after \passthrough{\lstinline!request.POST!}?}'' (Django ticket \href{https://code.djangoproject.com/ticket/27592}{27592}, \passthrough{\lstinline!pinned\_commit = django-4.2.11!}). Older Stack Overflow answers state Django \emph{always} raises \passthrough{\lstinline!RawPostDataException!} after \passthrough{\lstinline!POST!}; the Django 4.1 release notes clarify it is raised \emph{only for form-encoded} requests. The web is awash in the older ``always raises'' answers.
\item
  \textbf{\passthrough{\lstinline!requests\_retry\_adapter!}} --- ``\emph{Correct retry-strategy spelling for \passthrough{\lstinline!requests!} + urllib3 2.x}'' (\passthrough{\lstinline!pinned\_commit = requests-2.32.3+urllib3-2.2.2!}). Older Stack Overflow answers recommend \passthrough{\lstinline!urllib3.util.Retry(method\_whitelist=...)!}; urllib3 1.26 deprecated this in favour of \passthrough{\lstinline!allowed\_methods!} and urllib3 2.0 removed \passthrough{\lstinline!method\_whitelist!} entirely.
\item
  \textbf{\passthrough{\lstinline!pandas\_iterrows\_dtype!}} --- ``\emph{Does \passthrough{\lstinline!DataFrame.iterrows!} preserve column dtypes?}'' (\passthrough{\lstinline!pinned\_commit = pandas-2.2.2!}, \href{https://github.com/pandas-dev/pandas/issues/15014}{pandas\#15014}). Dozens of older answers say yes; the pandas 2.x reference says no --- rows are returned as \passthrough{\lstinline!Series!} and promoted to a \emph{common dtype} (typically \passthrough{\lstinline!object!}), with \passthrough{\lstinline!itertuples()!} recommended when dtypes must be preserved.
\end{enumerate}

For each issue we curated 3--5 evidence items (\passthrough{\lstinline!EvidenceItem!}s, see \passthrough{\lstinline!experiments/v2/p6b/case\_data.py!}) tagged with \passthrough{\lstinline!precision!}, \passthrough{\lstinline!condition!}, \passthrough{\lstinline!stale!}, and where appropriate \passthrough{\lstinline!superseded\_by\_id!}.

\subsection{Arms}\label{arms}

Both arms are deterministic and LLM-free; the difference is purely the \emph{discipline} applied:

\begin{itemize}
\item
  \textbf{\passthrough{\lstinline!without\_harness!} (baseline).} Process evidence in \emph{discovery order}; commit to the \emph{first-seen} qualifier as the final answer. This simulates the documented \emph{Lost-in-the-Middle} anchoring bias \citep{liu2023lostmiddle} where an unaided LLM agent latches onto the earliest-surfaced (and therefore typically most-popular-and-stale) result.
\item
  \textbf{\passthrough{\lstinline!with\_harness!} (treatment).} Process evidence in \emph{discovery order}; for each new item, call \passthrough{\lstinline!sublate\_with\_evidence!} on any pre-existing item it explicitly supersedes (via \passthrough{\lstinline!superseded\_by\_id!}) \emph{or} on any pre-existing stale item under the same \passthrough{\lstinline!(qualificand, condition)!} whose \passthrough{\lstinline!precision!} is dominated by the new item. Then run \passthrough{\lstinline!compact!} to drop sublated items from the live set. Final answer is the qualifier-set of the surviving live items.
\end{itemize}

\subsection{Headline numbers}\label{headline-numbers}

{\def\LTcaptype{none} % do not increment counter
\begin{longtable}[]{@{}
  >{\raggedright\arraybackslash}p{(\linewidth - 6\tabcolsep) * \real{0.2500}}
  >{\raggedright\arraybackslash}p{(\linewidth - 6\tabcolsep) * \real{0.2500}}
  >{\raggedright\arraybackslash}p{(\linewidth - 6\tabcolsep) * \real{0.2500}}
  >{\raggedright\arraybackslash}p{(\linewidth - 6\tabcolsep) * \real{0.2500}}@{}}
\toprule\noalign{}
\begin{minipage}[b]{\linewidth}\raggedright
metric
\end{minipage} & \begin{minipage}[b]{\linewidth}\raggedright
with-harness
\end{minipage} & \begin{minipage}[b]{\linewidth}\raggedright
without-harness
\end{minipage} & \begin{minipage}[b]{\linewidth}\raggedright
delta
\end{minipage} \\
\midrule\noalign{}
\endhead
\bottomrule\noalign{}
\endlastfoot
Cases correct & \textbf{3 / 3} (100\%) & 0 / 3 (0\%) & \textbf{+3} \\
Forbidden-claim hits (lower-is-better) & \textbf{1} & 5 & \textbf{−4} \\
Stale items in final live set (LIB) & \textbf{0} & 3 & \textbf{−3} \\
Total \passthrough{\lstinline!sublate\_with\_evidence!} events fired & \textbf{7} & 0 & --- \\
Total \passthrough{\lstinline!compact!} events fired & \textbf{3} & 0 & --- \\
Context tokens used (matched) & 213 / 128 / 138 & 213 / 128 / 138 & 0 \\
\end{longtable}
}

Source: \passthrough{\lstinline!experiments/results/p6b/\_summary.json!}. Visualised in Figures\textasciitilde{}\ref{fig:f08} (per-case accuracy) and \ref{fig:f09} (forbidden-claim hits); per-case detail is in Table\textasciitilde{}\ref{tab:t3_p6b}.

\begin{figure}[!tbp]
\centering
\begin{subfigure}[t]{0.48\linewidth}
  \centering
  \includegraphics[width=\linewidth]{F08_P6B_per_case_accuracy.png}
  \caption{Per-case correctness (3/3 vs.\ 0/3).}
  \label{fig:f08}
\end{subfigure}\hfill
\begin{subfigure}[t]{0.48\linewidth}
  \centering
  \includegraphics[width=\linewidth]{F09_P6B_forbidden_claims.png}
  \caption{Forbidden-claim hits (lower is better).}
  \label{fig:f09}
\end{subfigure}
\caption{\textbf{P6-B live case study.} Three real-world stale-vs-fresh evidence trails (Django, requests, pandas) processed under \emph{identical} token budgets. The system fires 7 \texttt{sublate\_with\_evidence} events and 3 \texttt{compact} events; the baseline anchors on the first-seen (and typically stale) item.}
\label{fig:f08f09}
\end{figure}

\begin{table}[!tbp]
\centering
\caption{\textbf{T3 --- P6-B per-case detail.} Sublations fired, compactions, forbidden-claim audit, and final-answer correctness for each of the three Python issues. Source: \texttt{experiments/results/p6b/\_summary.json}.}
\label{tab:t3_p6b}
{\scriptsize\setlength{\tabcolsep}{3pt}%
{\def\LTcaptype{none} % do not increment counter
\begin{longtable}[]{@{}
  >{\raggedright\arraybackslash}p{(\linewidth - 12\tabcolsep) * \real{0.1429}}
  >{\raggedright\arraybackslash}p{(\linewidth - 12\tabcolsep) * \real{0.1429}}
  >{\raggedright\arraybackslash}p{(\linewidth - 12\tabcolsep) * \real{0.1429}}
  >{\raggedright\arraybackslash}p{(\linewidth - 12\tabcolsep) * \real{0.1429}}
  >{\raggedright\arraybackslash}p{(\linewidth - 12\tabcolsep) * \real{0.1429}}
  >{\raggedright\arraybackslash}p{(\linewidth - 12\tabcolsep) * \real{0.1429}}
  >{\raggedright\arraybackslash}p{(\linewidth - 12\tabcolsep) * \real{0.1429}}@{}}
\toprule\noalign{}
\begin{minipage}[b]{\linewidth}\raggedright
case\_id
\end{minipage} & \begin{minipage}[b]{\linewidth}\raggedright
correct\_with
\end{minipage} & \begin{minipage}[b]{\linewidth}\raggedright
correct\_without
\end{minipage} & \begin{minipage}[b]{\linewidth}\raggedright
forbidden\_with
\end{minipage} & \begin{minipage}[b]{\linewidth}\raggedright
forbidden\_without
\end{minipage} & \begin{minipage}[b]{\linewidth}\raggedright
stale\_in\_ctx\_delta
\end{minipage} & \begin{minipage}[b]{\linewidth}\raggedright
ctx\_tokens
\end{minipage} \\
\midrule\noalign{}
\endhead
\bottomrule\noalign{}
\endlastfoot
django\_request\_body & 1 & 0 & 0 & 1 & -1 & 213 \\
requests\_retry\_adapter & 1 & 0 & 1 & 2 & -1 & 128 \\
pandas\_iterrows\_dtype & 1 & 0 & 0 & 2 & -1 & 138 \\
\end{longtable}
}
}
\end{table}

The \emph{context tokens used} row is critical: both arms operate under \emph{identical} token budgets. The harness wins not by buying more context but by \emph{using the same budget more disciplinedly}. The seven \passthrough{\lstinline!sublate\_with\_evidence!} events and three \passthrough{\lstinline!compact!} events are exactly the operational footprint of the Section 4.3 / 4.7 commitments.

\subsection{Per-case narrative}\label{per-case-narrative}

\subsubsection{\texorpdfstring{\texttt{django\_request\_body}}{django\_request\_body}}\label{django_request_body}

The baseline anchored on the pre-3.2 advice (the first item in discovery order, since it was the most-upvoted Stack Overflow answer when the issue was filed). The harness saw the same item, accepted it provisionally, then fired \passthrough{\lstinline!sublate\_with\_evidence!} against the stale Stack Overflow snippets as fresher items arrived (the official Django 4.1 release notes and the Django 4.2 docs); per \passthrough{\lstinline!experiments/results/p6b/django\_request\_body.json!} (\passthrough{\lstinline!metrics.sublations!}), this case fires \textbf{3 sublations} and \textbf{1 compaction}, after which the final live set contains only the fresh Django 4.x items (\passthrough{\lstinline!store\_size\_active = 4!}, \passthrough{\lstinline!store\_size\_sublated = 1!}). The final answer correctly cited the form-encoded clarification from the 4.1 release notes.

The forbidden-claim audit: the baseline emitted one forbidden phrase (\passthrough{\lstinline!Django 1.x!}/\passthrough{\lstinline!always raises!}/\passthrough{\lstinline!never raises!}); the harness emitted zero.

\subsubsection{\texorpdfstring{\texttt{requests\_retry\_adapter}}{requests\_retry\_adapter}}\label{requests_retry_adapter}

Stale Stack Overflow snippets used the old \passthrough{\lstinline!Retry(method\_whitelist=...)!} parameter; the fresh urllib3 v2 doc snippet uses \passthrough{\lstinline!allowed\_methods=...!}. Per the per-case JSON (\passthrough{\lstinline!requests\_retry\_adapter.json!}, \passthrough{\lstinline!metrics.sublations!}), the harness fires \textbf{2 sublations} and \textbf{1 compaction} and commits to \passthrough{\lstinline!allowed\_methods!}. The baseline locked onto \passthrough{\lstinline!method\_whitelist!} and never recovered.

The forbidden-claim audit: the baseline emitted two forbidden references (e.g.~\passthrough{\lstinline!method\_whitelist!}, \passthrough{\lstinline!requests.packages.urllib3!}); the harness emitted one --- a single \passthrough{\lstinline!method\_whitelist!} mention surviving inside the \emph{sublation event's} explanatory \passthrough{\lstinline!reason!} string for audit purposes only, with zero references in the final live qualifier set.

\subsubsection{\texorpdfstring{\texttt{pandas\_iterrows\_dtype}}{pandas\_iterrows\_dtype}}\label{pandas_iterrows_dtype}

Stale snippets all said ``iterrows preserves dtype''; the fresh pandas-2.x doc snippet says ``no, use \passthrough{\lstinline!itertuples!} (or accept that rows are promoted to a common dtype)''. Per \passthrough{\lstinline!pandas\_iterrows\_dtype.json!} (\passthrough{\lstinline!metrics.sublations!}), the harness fires \textbf{2 sublations} and \textbf{1 compaction}, leaving only the fresh pandas-2.x items in the live set, and commits to the \emph{common-dtype / \passthrough{\lstinline!itertuples!}} answer. The baseline committed to the wrong answer.

The forbidden-claim audit: the baseline emitted two forbidden phrases (\passthrough{\lstinline!preserves the column dtype!} and/or \passthrough{\lstinline!you can rely on integer!}); the harness emitted zero.

\subsection{Why this matters}\label{why-this-matters}

The L2 case study is the smallest possible \emph{ecologically valid} test of the harness's central claim: \emph{the discipline pays off even with no model-generation step, even under identical token budgets, on real-world stale-vs-fresh evidence trails}. The 3-of-3 result is small in \emph{n} but \emph{epistemically large}: it shows the mechanism is not benchmark-specific. The benchmark we constructed for L3 (Section 10) makes the same point at \(n=120\) with multi-seed multi-model paired statistics.

\subsection{Reproducibility}\label{reproducibility}

The case-study runner is \passthrough{\lstinline!experiments/v2/p6b/run\_case\_study.py!}. The case data is \passthrough{\lstinline!experiments/v2/p6b/case\_data.py!}. Both are deterministic and LLM-free. Re-running \passthrough{\lstinline!python -m experiments.v2.p6b.run\_case\_study!} re-emits \passthrough{\lstinline!experiments/results/p6b/*.json!} byte-identically.
